\documentclass{article}
\usepackage[margin=1in]{geometry}
\usepackage{amsmath}
\usepackage{relsize}

\DeclareMathOperator{\Var}{Var}
\DeclareMathOperator{\Cov}{Cov}

\newcommand{\bmu}{\boldsymbol{\mu}}
\newcommand{\bw}{\mathbf{w}}
\newcommand{\bSigma}{\boldsymbol{\Sigma}}
\newcommand{\bone}{\mathbf{1}}

\title{Closed-Form Solution of Markowitz-Mean-Variance Optimization}
\author{Puthi Uddam SEAN}
\date{\today}

\begin{document}

\maketitle{}

Suppose we have a portfolio ${p}$ consiting of ${2}$ assets ${A}$ and ${B}$, with respective weights ${w_A}$ and ${w_B}$.
The Variance of the portfolio ${p}$ is given by:
\begin{equation}
    \Var(p) = \Var(w_A A + w_B B) = \Cov(w_A A + w_B B, w_A A + w_B B)
\end{equation} 

Thanks to the bilinearity of covariance, we can expand the above equation as follows :

\begin{equation}
    \begin{aligned}
        \Var(p) &= w_A^2 \Var(A) + w_B^2 \Var(B) + 2 w_A w_B \Cov(A, B) \\ 
        &= w_A^2 \sigma^2_A + w_B^2 \sigma^2_B + 2 w_A w_B \sigma_A \sigma_B \rho_{AB} \text{ , where } \rho_{AB} \text{ is the correlation coefficient.}
    \end{aligned}
\end{equation}

Thanks to this equation, we can see that indeed 2 assets with negative correlation would reduce the variance and therefore, the volatility squared of our portfolio.

Now, let's consider the case of ${n}$ assets, with respective weights ${\bw = (w_1, w_2, \ldots, w_n)}$ and returns ${r_1, r_2, \ldots, r_n}$. We thus have the return of the portfolio ${r_p = \sum_{i=1}^{n} w_i r_i}$. Treating ${r_i}$ as random variables, we can deduce the following:
\begin{equation}
    \Var(r_p) = \Cov(r_p, r_p) = \Cov\left(\sum_{i=1}^{n} w_i r_i, \sum_{j=1}^{n} w_j r_j\right) = \sum_{i=1}^{n} \sum_{j=1}^{n} w_i w_j \Cov(r_i, r_j)
\end{equation}

On the other hand, let $\bSigma$ be a covariance matrix of the returns of the assets, with $\Sigma_{ij} = \Cov(r_i, r_j)$. We therefore have $(\bSigma \bw)_i = \mathlarger{\sum_{j=1}^n \Sigma_{ij} w_j}$ and consequently 
\begin{equation}
\bw^\top \bSigma \bw = \sum_{i=1}^{n} w_i (\bSigma \bw)_i = \sum_{i=1}^{n} w_i \sum_{j=1}^{n} \Sigma_{ij} w_j = \sum_{i=1}^{n} \sum_{j=1}^{n} w_i w_j \Cov(r_i, r_j) = \Var(r_p)
\end{equation}

\end{document}